# Capital Grow System

## Формальная математическая спецификация

## 1. Назначение

Capital Grow System — это алгоритмическая система управления капиталом, предназначенная для дискретного перераспределения активов портфеля в зависимости от положения цены BTC относительно фиксированной стартовой базы и заданных ценовых уровней.

Система предназначена для:

1. системного управления долями BTC, GM-позиции BTC/USD и USD;
2. наращивания BTC-эквивалента капитала на длинной дистанции;
3. увеличения долларовой стоимости портфеля;
4. ограничения риска ликвидации при использовании заёмных средств;
5. полного исключения субъективных решений в процессе работы.

---

## 2. Состав портфеля

Портфель состоит из следующих компонентов:

* (B_t) — количество BTC в момент времени (t);
* (G_t) — количество LP-единиц GM Pool BTC/USD в момент времени (t);
* (U_t) — количество USD-стейблкоинов в момент времени (t);
* (D_t) — объём долга в USD в момент времени (t).

Рыночные переменные:

* (P_t) — рыночная цена BTC в USD в момент времени (t);
* (M_t) — значение 200-дневной скользящей средней BTC в USD в момент времени (t);
* (Q_t) — рыночная цена одной LP-единицы GM Pool BTC/USD в USD в момент времени (t).

---

## 3. Стоимость портфеля

Чистая стоимость портфеля в USD определяется как:

[
V_t = B_t P_t + G_t Q_t + U_t - D_t
]

Все целевые доли активов в системе определяются относительно чистой стоимости (V_t).

---

## 4. Стоимость залога и риск

Стоимость залога определяется как:

[
C_t = B_t P_t
]

Кредитное плечо системы измеряется через отношение долга к залогу:

[
\mathrm{LTV}_t = \frac{D_t}{C_t}
]

При (C_t = 0) новые заёмные операции запрещены.

Если инфраструктура исполнения позволяет рассчитывать Health Factor, то дополнительно используется величина:

[
HF_t
]

с внешним определением по правилам кредитного протокола.

---

## 5. Параметры системы

### 5.1. Базовые параметры риска

Целевой максимум отношения долга к залогу:

[
L_{\mathrm{target}} = 0.20
]

Жёсткий аварийный порог отношения долга к залогу:

[
L_{\mathrm{hard}} = 0.25
]

Целевой минимальный Health Factor:

[
HF_{\mathrm{target}} = 1.80
]

Минимально допустимый Health Factor для открытия новых кредитных позиций:

[
HF_{\mathrm{min}} = 1.50
]

Размер одного кредитного шага:

[
\alpha = 0.03
]

Минимальная финальная доля BTC на верхней границе системы:

[
B_{\min} = 0.10
]

Минимальный относительный перевес нового ориентира над старым для перезапуска:

[
\varepsilon = 0.02
]

---

## 6. Инициализация системы

В момент запуска системы (t=t_0) фиксируются:

### 6.1. Стартовая цена

[
S = M_{t_0}
]

### 6.2. Стартовый ориентир BTC

[
O = \frac{V_{t_0}}{2S}
]

где

[
V_{t_0} = B_{t_0}P_{t_0} + G_{t_0}Q_{t_0} + U_{t_0} - D_{t_0}
]

Величина (O) — это стартовый ориентир BTC, выражающий количество BTC, соответствующее половине чистой стоимости портфеля по цене (S).

### 6.3. Служебные переменные состояния

Система хранит:

* (n_t) — индекс текущей активной ступени;
* (d_t \in {\uparrow,\downarrow,0}) — последнее подтверждённое направление движения;
* (s_t \in {0,1}) — флаг пропуска первой ступени после смены направления;
* (D_{\mathrm{peak}}) — максимальный достигнутый долг в текущем кредитном цикле.

На момент запуска:

[
d_{t_0}=0,\qquad s_{t_0}=0,\qquad D_{\mathrm{peak}}=D_{t_0}
]

---

## 7. Лестницы уровней

## 7.1. Верхняя лестница

Множество относительных уровней роста:

[
\mathcal{L}^{\uparrow} =
{0.12,0.24,0.36,0.48,0.60,0.72,0.84,0.96,1.08,1.20,1.32,1.44,1.56,1.68,1.80}
]

Абсолютные цены верхних ступеней:

[
P_k^{\uparrow} = S(1+\ell_k^{\uparrow}),
\qquad \ell_k^{\uparrow}\in \mathcal{L}^{\uparrow}
]

## 7.2. Нижняя лестница

Множество относительных уровней падения:

[
\mathcal{L}^{\downarrow} =
{0.11,0.21,0.30,0.38,0.45,0.51,0.56,0.60,0.63,0.65,0.66}
]

Абсолютные цены нижних ступеней:

[
P_k^{\downarrow} = S(1-\ell_k^{\downarrow}),
\qquad \ell_k^{\downarrow}\in \mathcal{L}^{\downarrow}
]

---

## 8. Целевые доли портфеля

Для каждой ступени задаются целевые доли:

* (w_B) — доля BTC;
* (w_G) — доля GM-позиции;
* (w_U) — доля USD.

Всегда выполняется:

[
w_B + w_G + w_U = 1
]

## 8.1. Целевые доли на верхней лестнице

| Уровень  | (w_B) | (w_G) | (w_U) |
| -------- | ----: | ----: | ----: |
| (S+12%)  |  0.40 |  0.60 |  0.00 |
| (S+24%)  |  0.30 |  0.70 |  0.00 |
| (S+36%)  |  0.20 |  0.80 |  0.00 |
| (S+48%)  |  0.10 |  0.90 |  0.00 |
| (S+60%)  |  0.00 |  1.00 |  0.00 |
| (S+72%)  |  0.00 |  0.90 |  0.10 |
| (S+84%)  |  0.00 |  0.80 |  0.20 |
| (S+96%)  |  0.00 |  0.70 |  0.30 |
| (S+108%) |  0.00 |  0.60 |  0.40 |
| (S+120%) |  0.10 |  0.40 |  0.50 |
| (S+132%) |  0.10 |  0.30 |  0.60 |
| (S+144%) |  0.10 |  0.20 |  0.70 |
| (S+156%) |  0.10 |  0.10 |  0.80 |
| (S+168%) |  0.10 |  0.05 |  0.85 |
| (S+180%) |  0.10 |  0.00 |  0.90 |

Для всех верхних ступеней начиная с (S+120%) должно выполняться ограничение:

[
w_B \ge B_{\min}
]

## 8.2. Целевые доли на нижней лестнице до кредитной зоны

| Уровень | (w_B) | (w_G) | (w_U) |
| ------- | ----: | ----: | ----: |
| (S-11%) |  0.60 |  0.40 |  0.00 |
| (S-21%) |  0.70 |  0.30 |  0.00 |
| (S-30%) |  0.80 |  0.20 |  0.00 |
| (S-38%) |  0.90 |  0.10 |  0.00 |
| (S-45%) |  1.00 |  0.00 |  0.00 |

## 8.3. Кредитная зона

На уровнях ниже (S-45%) базовые целевые доли фиксируются как:

[
w_B = 1,\qquad w_G = 0,\qquad w_U = 0
]

Кредитные ступени располагаются на уровнях:

* (S-51%),
* (S-56%),
* (S-60%),
* (S-63%),
* (S-65%),
* (S-66%).

На этих уровнях допускается открытие новых кредитных шагов в соответствии с разделом 12.

---

## 9. Начальное распределение капитала

## 9.1. Старт при совпадении цены и стартовой базы

Если

[
P_{t_0} = S
]

то стартовая структура портфеля должна быть приведена к долям:

[
w_B = 0.50,\qquad w_G = 0.50,\qquad w_U = 0
]

## 9.2. Старт при отклонении цены от стартовой базы

Если

[
P_{t_0} > S
]

то выбирается максимальная верхняя ступень (k), для которой:

[
P_k^{\uparrow} \le P_{t_0}
]

Если

[
P_{t_0} < S
]

то выбирается максимальная нижняя ступень (k), для которой:

[
P_k^{\downarrow} \ge P_{t_0}
]

После выбора ступени портфель приводится к её целевым долям.

---

## 10. Правило определения активной ступени

В любой момент времени активная ступень определяется как наиболее дальняя от стартовой цены ступень, уже пересечённая текущей ценой.

Если (P_t \ge S), то активной считается максимальная верхняя ступень (k), для которой:

[
P_t \ge P_k^{\uparrow}
]

Если (P_t \le S), то активной считается максимальная нижняя ступень (k), для которой:

[
P_t \le P_k^{\downarrow}
]

Если ни одна ступень не пересечена, активной считается базовая стартовая зона.

---

## 11. Правило смены направления

Система использует правило пропуска первой ступени после смены направления.

### 11.1. Смена направления вверх на вниз

Если ранее было подтверждено направление

[
d_t = \uparrow
]

и цена начала двигаться вниз, то первая пересечённая вниз ступень не вызывает ребалансировку. В этот момент устанавливается:

[
s_t = 1
]

Следующая пересечённая вниз ступень вызывает ребалансировку и одновременно устанавливает:

[
d_t = \downarrow,\qquad s_t = 0
]

### 11.2. Смена направления вниз на вверх

Если ранее было подтверждено направление

[
d_t = \downarrow
]

и цена начала двигаться вверх, то первая пересечённая вверх ступень не вызывает ребалансировку. В этот момент устанавливается:

[
s_t = 1
]

Следующая пересечённая вверх ступень вызывает ребалансировку и одновременно устанавливает:

[
d_t = \uparrow,\qquad s_t = 0
]

### 11.3. Отсутствие смены направления

Если движение продолжается в уже подтверждённом направлении, то каждая следующая ступень вызывает ребалансировку без пропуска.

---

## 12. Приведение портфеля к целевым долям

Для текущей цены (P_t), текущей цены LP-единицы (Q_t), чистой стоимости (V_t) и целевых долей ((w_B,w_G,w_U)) целевые абсолютные объёмы активов определяются как:

[
B_t^\star = \frac{w_B V_t}{P_t}
]

[
G_t^\star = \frac{w_G V_t}{Q_t}
]

[
U_t^\star = w_U V_t
]

Если текущая ступень не принадлежит кредитной зоне, то ребалансировка заключается в переходе:

[
(B_t,G_t,U_t)\to(B_t^\star,G_t^\star,U_t^\star)
]

при неизменном долге (D_t), за исключением обязательных операций по погашению долга на отскоке согласно разделу 13.

---

## 13. Кредитный блок

## 13.1. Условие допуска к кредитным операциям

Новая кредитная операция допускается только при одновременном выполнении всех условий:

[
P_t \le 0.49S
]

[
\mathrm{LTV}*t < L*{\mathrm{target}}
]

и, если доступен Health Factor,

[
HF_t > HF_{\mathrm{target}}
]

Если

[
\mathrm{LTV}*t \ge L*{\mathrm{hard}}
]

или, при наличии Health Factor,

[
HF_t \le HF_{\mathrm{min}}
]

то новые кредитные операции запрещены.

## 13.2. Размер нового кредитного шага

На каждой кредитной ступени допустимая сумма нового займа определяется как:

[
\Delta D_t = \min\left(\alpha C_t,; D_{\max}(t)-D_t\right)
]

где

[
D_{\max}(t) = L_{\mathrm{target}}, C_t
]

Если (\Delta D_t \le 0), кредитная операция не выполняется.

После получения займа все заёмные USD немедленно конвертируются в BTC:

[
\Delta B_t = \frac{\Delta D_t}{P_t}
]

После чего состояние обновляется:

[
D_t \leftarrow D_t + \Delta D_t
]

[
B_t \leftarrow B_t + \Delta B_t
]

[
D_{\mathrm{peak}} \leftarrow \max(D_{\mathrm{peak}}, D_t)
]

## 13.3. Ограничение числа кредитных шагов

Каждая кредитная ступень может активировать не более одного кредитного шага в пределах одного кредитного цикла.

Кредитный цикл начинается в момент первого открытия долга после входа цены в кредитную зону и завершается в момент полного погашения долга:

[
D_t = 0
]

---

## 14. Погашение долга на восстановлении цены

При движении цены вверх после наличия долга система обязана сокращать долг поэтапно.

Целевой долг определяется как функция максимального долга текущего кредитного цикла (D_{\mathrm{peak}}):

[
D_t^\star =
\begin{cases}
0.75D_{\mathrm{peak}}, & P_t \ge 0.55S \
0.50D_{\mathrm{peak}}, & P_t \ge 0.62S \
0.25D_{\mathrm{peak}}, & P_t \ge 0.70S \
0, & P_t \ge 0.79S
\end{cases}
]

При пересечении соответствующего уровня вверх выполняется погашение долга до целевого значения (D_t^\star).

Если

[
D_t > D_t^\star
]

то для погашения долга продаётся часть BTC:

[
\Delta B_{\mathrm{repay}} = \frac{D_t - D_t^\star}{P_t}
]

Затем состояние обновляется:

[
B_t \leftarrow B_t - \Delta B_{\mathrm{repay}}
]

[
D_t \leftarrow D_t^\star
]

После достижения

[
D_t = 0
]

кредитный цикл считается завершённым, а значение (D_{\mathrm{peak}}) обнуляется:

[
D_{\mathrm{peak}} \leftarrow 0
]

---

## 15. Полный алгоритм обработки нового рыночного наблюдения

Для каждого нового момента времени (t) выполняется следующая последовательность действий.

### Шаг 1. Обновление рыночных значений

Получить:

[
P_t,\quad M_t,\quad Q_t
]

### Шаг 2. Обновление чистой стоимости портфеля

[
V_t = B_tP_t + G_tQ_t + U_t - D_t
]

### Шаг 3. Определение активной ступени

Определить текущую активную ступень относительно стартовой цены (S).

### Шаг 4. Проверка пересечения новой ступени

Если новая ступень не пересечена, состояние портфеля не изменяется.

Если новая ступень пересечена, перейти к шагу 5.

### Шаг 5. Проверка смены направления

Если пересечение произошло в направлении, противоположном (d_t), и флаг пропуска не установлен, то:

[
s_t \leftarrow 1
]

ребалансировка не выполняется.

Если пересечение произошло в том же направлении, либо пропуск уже был использован, перейти к шагу 6.

### Шаг 6. Выполнение операции на ступени

#### 6.1. Если ступень обычная

Портфель приводится к целевым долям ступени по формулам раздела 12.

#### 6.2. Если ступень кредитная

Сначала выполняется приведение к базовой структуре кредитной зоны:

[
w_B = 1,\qquad w_G = 0,\qquad w_U = 0
]

Затем выполняется проверка условий кредитования из раздела 13.
Если условия выполнены, открывается новый кредитный шаг.

### Шаг 7. Погашение долга на отскоке

Если цена пересекла вверх один из уровней восстановления:

* (0.55S),
* (0.62S),
* (0.70S),
* (0.79S),

то долг сокращается до целевого значения по формулам раздела 14.

### Шаг 8. Обновление направления

После подтверждённой операции устанавливается:

* при движении вверх:
  [
  d_t = \uparrow,\qquad s_t = 0
  ]

* при движении вниз:
  [
  d_t = \downarrow,\qquad s_t = 0
  ]

---

## 16. Перезапуск системы

## 16.1. Условие рассмотрения перезапуска

Перезапуск системы рассматривается только в момент пересечения ценой BTC своей 200DMA.

Формально, в момент (t_r) условие пересечения задаётся как:

[
(P_{t_r^-} - M_{t_r^-})(P_{t_r} - M_{t_r}) \le 0
]

## 16.2. Новая стартовая цена

При рассмотрении перезапуска новая стартовая цена определяется как:

[
S_{\mathrm{new}} = M_{t_r}
]

## 16.3. Новая чистая стоимость

Чистая стоимость портфеля в момент рассмотрения перезапуска:

[
V_{t_r} = B_{t_r}P_{t_r} + G_{t_r}Q_{t_r} + U_{t_r} - D_{t_r}
]

## 16.4. Новый ориентир BTC

[
O_{\mathrm{new}} = \frac{V_{t_r}}{2S_{\mathrm{new}}}
]

## 16.5. Правило перезапуска

Перезапуск разрешён только если:

[
O_{\mathrm{new}} > O_{\mathrm{old}}(1+\varepsilon)
]

При выполнении условия перезапуска система обновляет параметры:

[
S \leftarrow S_{\mathrm{new}}
]

[
O \leftarrow O_{\mathrm{new}}
]

[
d_t \leftarrow 0
]

[
s_t \leftarrow 0
]

Затем портфель приводится к стартовой структуре, соответствующей текущему положению цены относительно новой стартовой цены, по правилам раздела 9.

Если условие не выполнено, система продолжает работать без перезапуска.

---

## 17. Учёт GM-позиции

GM-позиция учитывается исключительно по рыночной стоимости:

[
V_t^{(G)} = G_tQ_t
]

Система не использует приближённые модели постоянной BTC-экспозиции GM-позиции.
Все операции с GM-позицией осуществляются только через её текущую рыночную цену (Q_t).

---

## 18. Метрики состояния и эффективности

## 18.1. BTC-эквивалент чистой стоимости портфеля

[
X_t = \frac{V_t}{P_t}
]

Величина (X_t) показывает, сколько BTC можно приобрести на текущую чистую стоимость портфеля.

## 18.2. Превышение стартового ориентира BTC

[
R_t = \frac{X_t}{O}
]

Если

[
R_t > 1
]

то система превысила стартовый ориентир BTC.

## 18.3. Долларовая доходность

[
Y_t = \frac{V_t}{V_{t_0}} - 1
]

## 18.4. Максимальная просадка

[
DD_t = 1 - \frac{V_t}{\max_{\tau\le t}V_\tau}
]

---

## 19. Инварианты корректной работы системы

Во все моменты времени должны выполняться следующие условия:

1. неотрицательность активов:
   [
   B_t \ge 0,\qquad G_t \ge 0,\qquad U_t \ge 0
   ]

2. неотрицательность долга:
   [
   D_t \ge 0
   ]

3. сумма целевых долей равна единице:
   [
   w_B + w_G + w_U = 1
   ]

4. кредитные операции разрешены только в кредитной зоне:
   [
   P_t \le 0.49S
   ]

5. отношение долга к залогу для новых кредитных операций должно удовлетворять:
   [
   \mathrm{LTV}*t < L*{\mathrm{target}}
   ]

6. при достижении аварийного порога новые кредитные операции запрещены:
   [
   \mathrm{LTV}*t \ge L*{\mathrm{hard}}
   ]

7. если доступен Health Factor, то для новых кредитных операций должно выполняться:
   [
   HF_t > HF_{\mathrm{target}}
   ]

8. если доступен Health Factor и
   [
   HF_t \le HF_{\mathrm{min}}
   ]
   новые кредитные операции запрещены;

9. на верхних финальных ступенях сохраняется минимальная доля BTC:
   [
   w_B \ge B_{\min}
   ]

10. каждая кредитная ступень может быть использована не более одного раза в рамках одного кредитного цикла.

---

## 20. События, изменяющие состояние системы

Состояние системы изменяется только при наступлении одного из следующих событий:

1. запуск системы;
2. пересечение новой ценовой ступени;
3. подтверждённая смена направления после пропуска первой ступени;
4. открытие нового кредитного шага;
5. обязательное частичное погашение долга на восстановлении;
6. полный выход из долга;
7. перезапуск системы при пересечении 200DMA и выполнении условия на новый ориентир BTC.

Во всех остальных случаях система не выполняет операций.

---

## 21. Завершённое определение системы

Capital Grow System полностью определяется следующими объектами:

1. состоянием портфеля
   [
   (B_t,G_t,U_t,D_t)
   ]

2. рыночными переменными
   [
   (P_t,M_t,Q_t)
   ]

3. стартовой ценой
   [
   S
   ]

4. ориентиром BTC
   [
   O
   ]

5. верхней и нижней лестницами ценовых уровней;

6. таблицей целевых долей портфеля;

7. правилом пропуска первой ступени после смены направления;

8. правилами кредитования и ограничениями по риску;

9. правилами погашения долга на восстановлении;

10. правилом перезапуска при пересечении 200DMA.
